\section{项目概述}

\subsection{背景：文献调研中的三重断裂}

过去十年，学术文献的检索入口越来越顺滑，但“检索”与“理解”之间的缝隙并没有缩小。一个刚进入新方向的研究生，通常会在一天之内经历三次断裂。

\begin{enumerate}
  \item \textbf{列表与结构之间的断裂。} 搜索引擎给出的是按相关度排序的标题列表；引用数据库给出的是有向的引用边。前者丢失了论文之间的关系，后者只说明“谁引用过谁”，却说不清“为什么这两篇值得放在一起读”。
  \item \textbf{关系与证据之间的断裂。} 通用问答模型可以流畅地总结一篇论文，但它的结论往往没有落点：用户看到“A 是 B 的基线”，却无法点回原文，确认这句话究竟来自摘要、引言，还是模型根据常识补出来的。
  \item \textbf{一次检索与长期调研之间的断裂。} 大多数工具把一次查询当成一次会话的终点。用户今天查一个方向，明天补两篇新论文，之前的图、笔记与问题很难继续生长；每次重新搜索，都像重新开始。
\end{enumerate}

我们认为，文献调研真正缺的不是“更多结果”，而是一张\textbf{能被质疑、能被累积、能被追问}的关系地图。StellarLit 因此把“论文”当作星体，把“关系”当作引力线：输入一篇种子论文，沿着引用与语义关系长出一片星空；每条边都带原文证据，每次回答都能点回出处。

\subsection{项目定位：从检索列表到研究星图}

StellarLit 的定位可以拆成三个转变。

\begin{itemize}
  \item \textbf{从“列表”到“星图”。} 论文不是一个孤立的条目，而是某个社区、某个领域中的节点。系统用三层图谱把结构事实、语义推断与算法派生分开表达，让用户既能看见局部论文，也能拉远看见领域星云。
  \item \textbf{从“引用”到“关系 + 证据”。} 除引用外，系统判定基线、同期、思路类似、同任务异方法、扩展、结论矛盾六类语义关系；每一类都有判定标准、反例与硬约束，每条语义边必须携带至少一条能指回原文的证据。
  \item \textbf{从“一次性问答”到“可生长的研究台”。} 同一片星空可以反复并入新种子，老节点坐标保持稳定；问答历史、决策轨迹与图像证据都被保留，用户可以回到上次的研究现场继续追问。
\end{itemize}



\subsection{设计目标}

\begin{table}[H]
  \centering
  \footnotesize
  \caption{StellarLit 的四项设计目标与可观察行为}
  \label{tab:goals}
  \begin{tabularx}{\textwidth}{L{2.6cm}L{5.6cm}Y}
    \toprule
\textbf{目标} & \textbf{含义} & \textbf{用户能观察到的行为} \\
    \midrule
    可信 & 结论有出处，推断可复核 & 点开任意语义边，看到关系类型、置信度与原文片段；回答中的角标可定位到论文与文本块 \\
    可长大 & 调研过程可累积，而非一次一图 & 同一片星空可多次并入种子；老星体不漂移，历史星空可一键回到上次视角 \\
    可对话 & 不只展示图，还能围绕图提问 & 流式回答、比较类问题拆解、证据不足时自动补查，并展示“为什么停” \\
    可扩展 & 能力不只服务站内页面 & 9 个只读 MCP 工具，外部 Agent 可读取星空、检索原文、查证关系 \\
    \bottomrule
  \end{tabularx}
\end{table}

\subsection{功能概览}

\begin{table}[H]
  \centering
  \caption{核心能力与实现要点}
  \label{tab:features}
  \begin{tabularx}{\textwidth}{L{3.1cm}L{4.6cm}Y}
    \toprule
\textbf{能力} & \textbf{用户入口} & \textbf{实现要点} \\
    \midrule
    种子解析 & 标题 / DOI / arXiv ID & 永远返回候选列表，由用户确认；候选带置信度与推荐理由，不替用户做决定 \\
    星空生长 & 选择候选后自动构建 & 展开引文邻域，去重、建图、社区检测、服务端布局；老节点作为锚点参与增量松弛 \\
    关系判定 & 构建时自动执行，可在设置页调整 & 硬约束预过滤 + 多标签 rubric 判定 + 引用片段核验；每条语义边带证据 \\
    证据查看 & 点击星体或关系边 & 展示摘要、关系类型、置信度、原文片段；边有固定的视觉编码 \\
    文献问答 & 底部观星台输入问题，可附 4 张图 & 混合检索 + Agent 补查 + SSE 流式回答；角标回原文，图像证据可定位 \\
    历史与账号 & 左侧历史、账号菜单 & 星空、问答历史、模型配置、用量与 MCP 令牌按账号隔离 \\
    模型服务 & 设置页按账号绑定 & 多厂商 OpenAI 兼容接口；判定 / 问答 / 翻译 / 向量模型分别选择 \\
    MCP & 设置页生成个人令牌 & 9 个只读工具，stdio 与 HTTP 双入口，外部 Agent 即插即用 \\
    \bottomrule
  \end{tabularx}
\end{table}

\subsection{与赛题要求的对应}

赛题强调“前端 Web 应用为主要呈现形态，后端服务仅作为 AI 能力的接入层”，并要求覆盖大模型集成、Agent 工作流、RAG、多模态、AI 安全与 MCP 等能力。本项目没有另起一套独立后端，而是把 AI 接入、轻量持久化与协议入口收敛在 Next.js 的 Route Handlers 中；图谱渲染、交互与 AI 决策轨迹仍在 Web 应用一侧完成。

\begin{table}[H]
  \centering
  \scriptsize
  \caption{赛题能力项与本项目实现的对应关系}
  \label{tab:contest}
  \begin{tabularx}{\textwidth}{L{3.4cm}L{7.2cm}L{2.4cm}}
    \toprule
\textbf{赛题能力} & \textbf{本项目实现} & \textbf{详见} \\
    \midrule
    大模型集成 & 多厂商 OpenAI 兼容接口 + Anthropic 适配；结构化 JSON 判定、SSE 流式回答、超时 / 限流 / 余额不足的可读降级 & 第 4.2--4.3 节 \\
    Agent 工作流 & 规划 → 工具调用 → 证据覆盖判定 → 补查或停止；最多 2 轮 / 5 步，界面留下完整决策轨迹 & 第 4.6 节 \\
    RAG & 摘要与关系原文切块；BM25 + 向量双路召回；RRF 融合、图结构扩展、意图重排与证据阈值 & 第 4.5 节 \\
    多模态 & 最多 4 张图；视觉模型抽取 OCR、表格、图表序列、关键数值与归一化区域坐标 & 第 4.7 节 \\
    AI 安全 & 输入净化、资料与指令隔离、IP 限流与并发上限、每日 token 预算、输出兜底 & 第 4.8 节 \\
    MCP & 9 个只读工具，个人访问令牌，stdio 与 HTTP 双入口 & 第 5.6 节 \\
    前端工程 & Next.js 15 + React 19 + TypeScript strict；Canvas 2D 自写渲染管线；ESLint 强制依赖单向 & 第 3 章 \\
    \bottomrule
  \end{tabularx}
\end{table}
